\documentclass[11pt]{article}
\usepackage{amsmath,amssymb,amsthm}
\usepackage[margin=1in]{geometry}
\usepackage{hyperref}

\newtheorem{lemma}{Lemma}
\newtheorem{theorem}{Theorem}
\newcommand{\ex}{\operatorname{ex}}

\title{A Star-Matching Dichotomy for the Family Extremal Number,\\
       after Hunter's Counterexample}
\author{KJ\_C}
\date{}

\begin{document}

\maketitle

\begin{abstract}
This note records a small dichotomy concerning Erd\H{o}s Problem~\#180 in
the regime where every individual forbidden graph has linear extremal
number. After Hunter's folklore counterexample to the original conjecture,
the open question is what remains true in this regime; the dichotomy
below is one such statement. The argument is elementary, the result is
likely folklore. The accompanying repository contains a Lean~4
formalization in which the corresponding theorem is proved from Lean's
foundational axioms only.
\end{abstract}

All graphs are finite and simple, and subgraph means not necessarily induced.
For a graph \(H\), write \(H^\circ\) for the graph obtained by deleting all
isolated vertices.

\section*{1. Background and statement}

Erd\H{o}s Problem~\#180 asks, in its original form, whether the family
extremal number \(\ex(n;\mathcal F)\) is determined by the individual
extremal numbers \(\ex(n;H)\) for \(H\in\mathcal F\), in some natural sense.
On the problem's forum, Zach Hunter recorded the following \emph{folklore
counterexample}: if \(\mathcal F=\{K_{1,a},\,bK_2\}\) with \(a,b\ge 2\),
then \(\ex(n;\mathcal F)=O_{\mathcal F}(1)\), while each individual
\(\ex(n;H_i)=\Theta(n)\). The conjecture as stated therefore fails in
general.

The remaining open question, framed on the problem page, is what holds
\emph{outside} this obstruction. The dichotomy below addresses this in the
all-linear regime: when every \(\ex(n;H)\) for \(H\in\mathcal F\) is
\(\Theta(n)\), the family extremal number is either \(\Theta(1)\) (and
this happens precisely when the star/matching obstruction is present) or
\(\Theta(n)\) (otherwise).

We do not claim the dichotomy itself is new. The proof uses only standard
tools and the result is plausibly folklore; the Lean formalization in the
accompanying repository is the contribution this note primarily exists to
document.

\section*{2. The single-forbidden-graph lemma (used as input)}

The proof relies on the following standard characterization, used as a
black box.

\begin{lemma}[Used; not proved here]
Let \(H\) be a finite graph. Then \(\ex(n;H)=\Theta(n)\) if and only if
\(H^\circ\) is a forest with at least two edges.
\end{lemma}

This is folklore in extremal graph theory; references include
Bollob\'as's \emph{Extremal Graph Theory} and Diestel's \emph{Graph
Theory}. The forward direction (forest with \(\ge 2\) edges \(\Rightarrow\)
\(\Theta(n)\)) follows from a degeneracy argument paired with star or
matching constructions. The converse (\(H^\circ\) contains a cycle
\(\Rightarrow\) \(\ex(n;H)\) is super-linear) follows from
Erd\H{o}s's high-girth/high-chromatic-number probabilistic construction.

In the Lean formalization, only the implication
\[
\ex(n;H)=\Theta(n) \;\Longrightarrow\; H^\circ \text{ has at least two edges}
\]
is actually used by the dispatch in Theorem~\ref{thm:main} below; the
other direction and the forest conclusion are not consumed. The
formalization therefore proves a weakened version sufficient for the
dispatch.

\section*{3. The dichotomy}

\begin{theorem}\label{thm:main}
Let \(\mathcal F\) be a nonempty finite family of finite graphs with
\[
\ex(n;H)=\Theta(n) \qquad \text{for every } H\in\mathcal F.
\]
Then exactly one of the following holds:

\begin{enumerate}
\item[\textup{(i)}] \(\mathcal F\) contains, after deleting isolated vertices,
both \(K_{1,a}\) and \(bK_2\) for some \(a,b\ge 2\). In this case
\(\ex(n;\mathcal F)=\Theta(1)\).

\item[\textup{(ii)}] \(\mathcal F\) contains no such star/matching pair. In
this case \(\ex(n;\mathcal F)=\Theta(n)\).
\end{enumerate}
\end{theorem}

\begin{proof}
By the lemma above, every \(H^\circ\) with \(H\in\mathcal F\) is a forest
with at least two edges.

\emph{Case (i).} Suppose \(\mathcal F\) contains both \(K_{1,a}\) and
\(bK_2\) with \(a,b\ge 2\). Let \(G\) be \(\mathcal F\)-free. Then
\(\Delta(G)\le a-1\) (else \(G\) contains \(K_{1,a}\)) and
\(\nu(G)\le b-1\) (else \(G\) contains \(bK_2\)). Let \(M\) be a maximum
matching. Since \(M\) is maximal, every edge of \(G\) has an endpoint in
\(V(M)\). Therefore
\[
e(G)\le\sum_{v\in V(M)}d_G(v)\le |V(M)|\cdot\Delta(G)\le 2(b-1)(a-1).
\]
So \(\ex(n;\mathcal F)=O(1)\). Since every member of \(\mathcal F\) has
at least two edges, a single edge is \(\mathcal F\)-free, giving
\(\ex(n;\mathcal F)\ge 1\). Hence \(\ex(n;\mathcal F)=\Theta(1)\).

\emph{Case (ii).} Suppose \(\mathcal F\) contains no such star/matching
pair. The upper bound is immediate: choose any \(H_0\in\mathcal F\); every
\(\mathcal F\)-free graph is \(H_0\)-free, so
\(\ex(n;\mathcal F)\le\ex(n;H_0)=O(n)\).

For the lower bound: if no member of \(\mathcal F\) is, after deleting
isolated vertices, a star, then \(K_{1,n-1}\) is \(\mathcal F\)-free
(every nonempty subgraph of a star is a star after deleting isolated
vertices), giving \(\ex(n;\mathcal F)\ge n-1\). Otherwise, some member
is a star, but by hypothesis no member is a matching with at least two
edges; hence a maximum matching on \(n\) vertices is \(\mathcal F\)-free,
giving \(\ex(n;\mathcal F)\ge\lfloor n/2\rfloor\). In both subcases
\(\ex(n;\mathcal F)=\Omega(n)\), and combined with the upper bound we
obtain \(\ex(n;\mathcal F)=\Theta(n)\).
\end{proof}

\section*{4. Notes}

\paragraph{Relation to \#180.} Theorem~\ref{thm:main} is not a resolution
of \#180. It addresses the all-linear regime under the assumption that
every individual member has linear extremal number, and within that
regime exhibits the star/matching obstruction (Hunter's case) as the
unique pattern producing constant family extremal number.

\paragraph{Constants in Case~(i).} The argument above gives the explicit
bound \(e(G)\le 2(a-1)(b-1)\). The Lean formalization establishes a
slightly weaker bound, \(O(|V(F_{\mathrm{star}})|\cdot |V(F_{\mathrm{matching}})|)\),
because of how isolated vertices in family members were handled in the
formalization; both bounds give \(\Theta(1)\), but the explicit constants
differ.

\paragraph{What the formalization establishes.} The Lean theorem
\verb|familiesTheorem| is exactly the statement of Theorem~\ref{thm:main}
above, formalized in Lean~4 against \texttt{Mathlib4}. After replacement
of all project-specific axioms,
\verb|#print axioms familiesTheorem| reports only Lean's foundational
axioms (\texttt{propext}, \texttt{Classical.choice}, \texttt{Quot.sound}).
The formalization was carried out in approximately eight hours of
agent wall-clock time.

\paragraph{Acknowledgments.} The folklore counterexample is due to Zach
Hunter, recorded on the Erd\H{o}s Problems forum thread for \#180. The
problem itself is due to Erd\H{o}s and Simonovits. The accompanying Lean
formalization was carried out by AI agents (Claude Sonnet~4.6 for initial
problem identification, Claude Opus~4.7 for prompt design and
verification, GPT-5.5 xhigh for Lean code generation).

\paragraph{Document provenance.} This note was drafted by Claude
Opus~4.7.

\end{document}
